% GENERATED manually from Kam (2009), Party Discipline and Parliamentary Politics, Ch. 3.
% Do not edit by hand.
\begin{table}[htbp]\centering
\caption{Backbench dissent in four Westminster systems, 1945--2005 (Kam 2009)}
\begin{threeparttable}
\begin{tabular}{lp{9.5cm}}
\toprule
Country & Pattern \\
\midrule
UK & Frequency $\sim$5\% pre-1970 $\to$ surged to $\sim$20\% (Heath 1970--4, Wilson/Callaghan 1974--9) $\to$ settled 10--15\% (1980s--90s) $\to$ 21\% (Labour, 2001--5). A dated structural break, c.\ 1970. \\
Canada & Similar \emph{average} level to the UK, but far more volatile -- swings between $>$50\% dissenting divisions in one term and near-zero the next; driven by minority-government dynamics and chronic mavericks rather than sustained party-wide unrest. \\
Australia & Much lower than UK/Canada. House rarely exceeds 5\%; the Labor Party was essentially perfectly unified in the House (1 dissenting vote / 3{,}020 sampled divisions). Dissent concentrates in the Senate instead, where the confidence convention does not bind. \\
New Zealand & Also low ($<$5\% dissenting divisions), though depth/extent grew from the mid-1980s and defections rose sharply after 1988, accelerating with the 1993 shift to PR. \\
\bottomrule
\end{tabular}
\begin{tablenotes}\footnotesize
\item Source: Kam (2009), \emph{Party Discipline and Parliamentary Politics}, Ch. 3, pp.\ 41--51.
\end{tablenotes}
\end{threeparttable}
\label{tab:ch02rebellioncomparative}
\end{table}
